\section{Introduction}
\label{sec:introduction}

The lithium battery industry spends billions of dollars each year pursuing ultra-high purity electrode materials and electrolytes. The prevailing wisdom is simple: purer is better. Yet recent experimental evidence tells a more nuanced story. \citet{choe2024reevaluation} demonstrated that ${\sim}1\%$ magnesium ``impurity'' in lithium carbonate precursors actually \emph{improved} cathode capacity retention by 6.5\%, while simultaneously reducing capital expenditure by 19.4\%. \citet{hobold2024lithiumoxide} showed that Li$_2$O---a species readily formed from trace water and oxygen---correlates more strongly with high Coulombic Efficiency (CE) than the widely targeted LiF. These findings raise a provocative question: {\bf can some trace impurities be deliberately engineered to improve battery performance?}

Understanding the role of trace impurities is particularly critical at the high-CE frontier. As \citet{hobold2021beyond} pointed out, when CE exceeds 99\%, trace impurities become ``inadvertent differentiators across electrolyte suppliers.'' At this regime, small differences in solid-electrolyte interphase (SEI) composition---driven by ppm-level contaminants---can determine whether a cell achieves 500 or 1{,}000 charge--discharge cycles. Despite this importance, no systematic study has mapped the dose-response relationships of multiple impurity types across the ppm concentration range relevant to manufacturing.

Prior work has studied individual impurities in isolation. \citet{fink2021metallic} examined metallic contaminants at high concentrations (1~wt\%), and \citet{baakes2023impact} modeled the safety implications of water impurities. However, these studies did not explore the beneficial regime at low concentrations or compare multiple impurity types within a unified framework. The gap is clear: we lack quantitative dose-response curves that predict which impurities help, which hurt, and at what concentrations the crossover occurs.

We address this gap through a three-pronged computational study. First, we use physics-based battery simulations (\pybamm Single Particle Model with electrolyte, SPMe) to model how four representative impurities---magnesium, water, iron, and oxygen-containing species---modify SEI growth kinetics and CE over 100 charge--discharge cycles. Second, we analyze the published Baakes et al.\ thermal abuse dataset to extract relationships between SEI composition and stability. Third, we fit mechanistic dose-response models (hormesis and monotonic) to our simulation and literature data to predict optimal impurity concentrations. Our key finding is that magnesium at ${\sim}1{,}000$~\ppm improves CE by $+0.0018\%$ ($p=0.018$) and capacity retention by $+0.10\%$, while iron at the same concentration degrades capacity retention by $-0.33\%$. SEI resistivity is the strongest predictor of impurity benefit ($r=-0.891$, $p<0.0001$).

We make the following contributions:
\begin{itemize}[leftmargin=*,itemsep=0pt,topsep=0pt]
    \item We conduct the first systematic computational study of trace impurity effects on CE across four impurity types (Mg, H$_2$O, Fe, O$_2$ species) at concentrations from 50 to 1{,}000~\ppm, identifying both beneficial and detrimental regimes.
    \item We establish SEI resistivity as the dominant mechanistic predictor of impurity benefit ($r=-0.891$), providing a simple design rule for impurity engineering: impurities that reduce SEI ionic resistance improve CE.
    \item We fit quantitative dose-response models and validate our predictions against published experimental data from four independent studies, finding consistent agreement on optimal concentrations and qualitative trends.
\end{itemize}

The remainder of this paper is organized as follows. \Secref{sec:related} reviews prior work on impurity effects and SEI engineering. \Secref{sec:method} describes our simulation methodology and dose-response modeling. \Secref{sec:results} presents simulation results and statistical analysis. \Secref{sec:discussion} discusses implications, limitations, and future directions. \Secref{sec:conclusion} concludes.
